%%
%% dk2nu.tex -- reference documentation for the dk2nu package
%%
%% Build (from this directory):   latexmk -pdf dk2nu.tex
%%   or:                          pdflatex dk2nu.tex ; pdflatex dk2nu.tex
%%
%% Code listings are pulled directly from the repository files by
%% relative path (../tree, ../genie, ../scripts, ../etc, ./examples),
%% so rebuilding the PDF picks up the current code.
%%
\documentclass[11pt]{article}

\usepackage[letterpaper,margin=1in]{geometry}
\usepackage[T1]{fontenc}
\usepackage{lmodern}
\usepackage{textcomp}
\usepackage{amsmath}
\usepackage{xspace}
\usepackage{booktabs}
\usepackage{longtable}
\usepackage{tabularx}
\usepackage{array}
\usepackage{enumitem}
\usepackage[table]{xcolor}
\usepackage{listings}
\usepackage[hidelinks]{hyperref}
\hypersetup{colorlinks=true,linkcolor=blue!50!black,urlcolor=blue!60!black,
            citecolor=black}

\setcounter{tocdepth}{2}
\setlength{\emergencystretch}{3em}

% ---------------------------------------------------------------------------
% listings
\definecolor{lstbg}{gray}{0.96}
\definecolor{lstcomment}{rgb}{0,0.45,0}
\definecolor{lstkeyword}{rgb}{0,0,0.6}
\definecolor{lststring}{rgb}{0.6,0.1,0.1}
\lstset{
  basicstyle=\ttfamily\scriptsize,
  backgroundcolor=\color{lstbg},
  commentstyle=\color{lstcomment},
  keywordstyle=\color{lstkeyword},
  stringstyle=\color{lststring},
  numbers=left, numberstyle=\tiny\color{gray}, stepnumber=1, numbersep=6pt,
  tabsize=2, showstringspaces=false, breaklines=true,
  columns=fullflexible, keepspaces=true,
  frame=single, rulecolor=\color{gray!50},
  upquote=true
}
\lstdefinestyle{cxx}{language=C++}
\lstdefinestyle{xml}{language=XML,
  % only "..." are strings, so apostrophes in comments (isn't) are harmless
  deletestring=[b]', deletestring=[d]',
  morekeywords={param_set,name,verbose,units,lunits,tunits,beamdir,type,
                rotation,axis,beampos,window,point,coord,enumax,reuse,
                upstreamz,minmaxwgt,maxwgtfail,partialread,using_param_set}}
\lstdefinestyle{plain}{language={},numbers=none}

% ---------------------------------------------------------------------------
% shorthands
\newcommand{\code}[1]{\texttt{#1}}
\newcommand{\genie}{\textsc{genie}\xspace}
\newcommand{\geant}{\textsc{geant}\xspace}
\newcommand{\ROOT}{\textsc{root}\xspace}
\newcommand{\fig}[1]{Figure~\ref{fig:#1}}
\newcommand{\dktwonu}{\texttt{dk2nu}\xspace}
\newcommand{\numu}{\ensuremath{\nu_{\mu}}\xspace}
\newcommand{\numubar}{\ensuremath{\bar{\nu}_{\mu}}\xspace}
\newcommand{\nue}{\ensuremath{\nu_{e}}\xspace}
\newcommand{\nuebar}{\ensuremath{\bar{\nu}_{e}}\xspace}
\newcommand{\tab}[1]{Table~\ref{tab:#1}}
\newcommand{\sect}[1]{Section~\ref{sec:#1}}
\newcommand{\app}[1]{Appendix~\ref{app:#1}}
\newcommand{\open}[1]{\textcolor{red!70!black}{\textbf{[open: #1]}}}

% field tables:  name | type | units | description
\newcolumntype{F}{>{\ttfamily\raggedright\arraybackslash}p{0.19\textwidth}}
\newcolumntype{T}{>{\ttfamily\raggedright\arraybackslash}p{0.13\textwidth}}
% narrower member, wider type columns (used by the Dk2Nu table)
\newcolumntype{G}{>{\ttfamily\raggedright\arraybackslash}p{0.15\textwidth}}
\newcolumntype{V}{>{\ttfamily\raggedright\arraybackslash}p{0.19\textwidth}}
\newcolumntype{U}{>{\raggedright\arraybackslash}p{0.07\textwidth}}
\newcolumntype{D}{>{\raggedright\arraybackslash}X}

% optional first argument: column spec (default "F T U D")
\newenvironment{fieldtable}[2][F T U D]
  {\par\noindent\small\tabularx{\textwidth}{#1}
   \toprule \textbf{\rmfamily member} & \textbf{\rmfamily type} &
   \textbf{units} & \textbf{description (#2)} \\ \midrule}
  {\bottomrule\endtabularx\par}

\title{The \dktwonu Neutrino Beam Simulation\\Flux Ntuple Format\\
       and \genie Flux Driver}
\author{Robert Hatcher\\Fermilab}
\date{Reference for \dktwonu v01\_12\_00 and later\\
      (\code{DK2NUVER 12}, \code{DKMETAVER 6})\\[1ex] \today}

\begin{document}
\maketitle

\begin{abstract}
\noindent
\dktwonu is a common \ROOT \code{TTree} format for the output of neutrino
beamline simulations (\code{g4numi}, \code{g4lbne}/\code{g4lbnf},
\code{flugg}, and others), together with code to fill it, to compute
neutrino energies and weights for arbitrary detector positions, to read it
from \ROOT, C++ and Python, and a \genie flux driver
(\code{genie::flux::GDk2NuFlux}) that turns the stored decays into neutrino
rays for event generation.  This document describes the package as it now
exists.  It replaces the 2012 proposal ``Proposal for a Unified `Flux'
N-tuple Format'' (MINOS-DocDB-9070), whose tables of proposed elements are
now superseded by the class definitions described here.
\end{abstract}

\tableofcontents
\clearpage

%===========================================================================
\section{Introduction}\label{sec:intro}

Accelerator neutrino beams are produced by protons striking a target; the
resulting hadrons are focused by magnetic horns and decay in flight in a
decay volume, producing neutrinos.  A beamline simulation records, for every
decay that produced a neutrino, enough information about the decay to
re-evaluate the probability (and the energy) of the neutrino passing through
\emph{any} point downstream.  The resulting ``flux'' files are therefore
really files of \emph{decays}, each with an importance weight, rather than
files of neutrinos.

Historically each beamline simulation (\code{gnumi} based on \geant{}3,
\code{flugg} based on \textsc{fluka}+\geant{}4, \code{g4numi}, the
\textsc{miner}$\nu$\textsc{a} variant of \code{g4numi}, \code{g4lbne})
wrote its own ntuple layout.  The layouts were similar but incompatible:
different basic types, capitalisation of leaf names, fixed array sizes, and
ad hoc additions.  The \dktwonu format replaces all of these with a single
layout that is
\begin{itemize}[nosep]
\item hierarchical (decay, neutrino rays, ancestry, target exit, \ldots);
\item variable length where appropriate (\code{std::vector} rather than fixed
      arrays, so no artificial overflow);
\item extensible without changing the format (user \code{vint}/\code{vdbl}
      vectors with names recorded in the metadata);
\item accompanied by a per-file metadata tree recording the protons-on-target
      (POT) and the exact detector locations used for any pre-computed rays.
\end{itemize}

\subsection{Package layout}\label{sec:layout}

\begin{center}\small
\begin{tabularx}{\textwidth}{>{\ttfamily}l X}
\toprule
\textrm{\textbf{directory}} & \textbf{contents} \\ \midrule
tree/    & the data classes (\code{bsim::Dk2Nu}, \code{bsim::DkMeta},
           \code{bsim::NuChoice}, \ldots) and helper functions
           (\code{readWeightLocations}, \code{calcLocationWeights},
           \code{calcEnuWgt}); built into \code{libdk2nuTree}. \\
genie/   & the \genie flux driver \code{genie::flux::GDk2NuFlux}; built into
           \code{libdk2nuGenie}. \\
python/  & \code{pybind11} bindings (\code{pydk2nu}) and examples. \\
scripts/ & \code{load\_dk2nu.C} for interactive \ROOT; \code{examples/}
           for filling/reading; \code{convert/} for converting legacy
           ntuples. \\
etc/     & text files of standard detector locations. \\
doc/     & this document. \\
cmake/   & platform specific CMake fragments. \\
\bottomrule
\end{tabularx}
\end{center}

All data classes and helper functions live in the C++ namespace
\code{bsim} (``beam simulation'').

\subsection{Requirements}
The package requires \ROOT 6 and C++ compiler support matching that \ROOT
build.  The \genie interface requires \genie v3 (R-3\_xx\_yy; the driver has
been updated for the 3.04 series), together with \code{libxml2} and
\code{log4cpp}.  Support for \ROOT 5 and \genie v2 has been dropped.  The
Python bindings additionally need Python~3 and CMake $\geq 3.14$.

%===========================================================================
\section{Conventions}\label{sec:conventions}

\paragraph{Units.}  Unless stated otherwise, all quantities stored in the
\dktwonu and \code{dkmeta} trees use
\begin{itemize}[nosep]
\item length: centimetres (cm);
\item energy and momentum: GeV and GeV/$c$;
\item time: nanoseconds (ns).
\end{itemize}
Some legacy simulations used mm and MeV internally; the converters
(\app{convert}) rescale on input.

\paragraph{Beam coordinate system.}  Positions and momenta are in the
\emph{beam} coordinate system of the simulation: $\hat z$ is along the
nominal beam (decay-pipe) axis, and $\hat y$ is vertical (up).  The origin
is whatever the beamline simulation uses for its own geometry; it must be
the same origin used for the locations recorded in the metadata
(\sect{dkmeta}) and for the beam-to-detector transformation given to \genie
(\sect{genie-coords}).

\paragraph{Particle codes.}  All particle species (\code{ntype},
\code{ptype}, \code{tptype}, \code{Ancestor::pdg}, \code{nucleus}) are
\textsc{pdg} codes.  (The legacy formats used \geant{}3 codes, with the
extension $\nu_e=53$, $\bar\nu_e=52$, $\nu_\mu=56$, $\bar\nu_\mu=55$;
the converters translate these.)

\paragraph{Default (``unset'') values.}  Every class provides
\code{clear()}, which resets its members to recognisably invalid values
(\code{tree/dflt.h}):
\code{bsim::kDfltInt} $=-9999$, \code{bsim::kDfltDouble} $=-9999.99$,
\code{bsim::kDfltFloat} $=-9999.99$, \code{bsim::kDfltUInt} $=9999$ and
\code{bsim::kDfltBool} $=$ \code{false}.  The overloaded function
\code{bsim::IsDefault(value)} tests for these.  Vectors are cleared to
zero length.  A few members use other sentinels; these are noted in the
tables below.

\paragraph{Class versions.}  All tree classes in \code{dk2nu.h} carry the
\ROOT class version \code{DK2NUVER} (currently 12), the metadata classes
in \code{dkmeta.h} carry \code{DKMETAVER} (currently 6) and
\code{NuChoice} carries version 2.  These must be incremented whenever a
data member changes.  The \ROOT dictionaries use schema evolution
(``\code{+}'' in \code{LinkDef.h}), so older files remain readable.

\paragraph{Common methods.}  Every class provides \code{clear()} and
\code{AsString()} (each taking an optional \code{std::string} option),
the latter returning a formatted description, and a global \code{operator<<}.  \code{bsim::Dk2Nu} also has
\code{Print()}.

%===========================================================================
\section{File structure}\label{sec:file}

A \dktwonu file contains two trees:

\begin{center}\small
\begin{tabular}{llll}
\toprule
\textbf{tree} & \textbf{branch} & \textbf{class} & \textbf{entries} \\
\midrule
\code{dk2nuTree}  & \code{dk2nu}  & \code{bsim::Dk2Nu}  & one per recorded decay \\
\code{dkmetaTree} & \code{dkmeta} & \code{bsim::DkMeta} & one per job (normally one per file) \\
\bottomrule
\end{tabular}
\end{center}

The metadata is kept as a tree, not as a single object, so that files can
be chained with \code{TChain} or merged with \code{hadd}; a merged file
simply has several \code{dkmeta} entries.  Entries in the two trees are
associated through the \code{job} number: a \code{Dk2Nu} entry belongs to
the \code{DkMeta} entry with the same \code{job}.  Job numbers should
therefore be unique across any set of files that will be used together.

The same incident proton can produce more than one entry, so
\code{potnum} may repeat; \code{(job, potnum, jobindx)} identifies an entry.
The number of protons-on-target represented by a file is \emph{only}
available from \code{DkMeta::pots}; it cannot be reliably inferred from the
range of \code{potnum}.

Trees are written with split level 99 so that individual members can be
read, drawn, or disabled independently (e.g.\ \code{decay.ntype}).
\fig{classes} summarises the containment.

\begin{figure}[htb]
\centering\small
\begin{tabular}{l}
\code{bsim::Dk2Nu} \\
\quad\code{job, potnum, jobindx} \\
\quad\code{decay} \hfill \code{bsim::Decay} \\
\quad\code{nuray[]} \hfill \code{std::vector<bsim::NuRay>} \\
\quad\code{ancestor[]} \hfill \code{std::vector<bsim::Ancestor>} \\
\quad\code{ppvx, ppvy, ppvz} \\
\quad\code{tgtexit} \hfill \code{bsim::TgtExit} \\
\quad\code{traj[]} \hfill \code{std::vector<bsim::Traj>} \\
\quad\code{flagbits, vint[], vdbl[]} \\[1ex]
\code{bsim::DkMeta} \\
\quad\code{job, pots} \\
\quad\code{beamsim, physics, physcuts, tgtcfg, horncfg, dkvolcfg} \\
\quad\code{beam0x, beam0y, beam0z, beamhwidth, beamvwidth, beamdxdz, beamdydz} \\
\quad\code{location[]} \hfill \code{std::vector<bsim::Location>} \\
\quad\code{vintnames[], vdblnames[]} \\
\end{tabular}
\caption{Containment of the \dktwonu and \code{dkmeta} classes.}
\label{fig:classes}
\end{figure}

%===========================================================================
\section{The decay tree: \code{bsim::Dk2Nu}}\label{sec:dk2nu}

\subsection{\code{bsim::Dk2Nu}}

{\setlength{\tabcolsep}{3pt}%
\begin{fieldtable}[G V U D]{Dk2Nu}
job      & Int\_t   & -- & job number; matches \code{DkMeta::job}. \\
potnum   & Int\_t   & -- & number of the proton (within the job) that
                            initiated this decay chain.  Cleared to $-1$. \\
jobindx  & Int\_t   & -- & index of the entry within the job; unique where
                            \code{potnum} is not.  Cleared to $-1$. \\
decay    & Decay    & -- & core decay information (\sect{decay}). \\
nuray    & vector<\hspace{0pt}NuRay> & -- & neutrino rays through the metadata
                            locations (\sect{nuray}). \\
ancestor & vector<\hspace{0pt}Ancestor> & -- & chain from primary proton to neutrino
                            (\sect{ancestor}). \\
ppvx, ppvy, ppvz & Double\_t & cm & production vertex of the neutrino
                            parent.  Redundant with the ancestor list
                            (\code{ancestor[indxp()].start\{x,y,z\}}) when
                            that is filled; kept for convenience and for
                            files without ancestry. \\
tgtexit  & TgtExit  & -- & ancestor as it left the target (\sect{tgtexit}). \\
traj     & vector<\hspace{0pt}Traj> & -- & optional trajectory points (\sect{traj}). \\
flagbits & Int\_t   & -- & bit flags (\sect{extend}).  Cleared to 0. \\
vint     & vector<\hspace{0pt}Int\_t>    & -- & user-defined integers (\sect{extend}). \\
vdbl     & vector<\hspace{0pt}Double\_t> & -- & user-defined doubles (\sect{extend}). \\
\end{fieldtable}}

\noindent Methods:
\begin{itemize}[nosep]
\item \code{size\_t indxnu() const} --- index of the neutrino in
      \code{ancestor} (\code{ancestor.size()-1});
\item \code{size\_t indxp() const} --- index of the neutrino's parent
      (\code{size()-2});
\item \code{size\_t indxgp() const} --- index of the grandparent
      (\code{size()-3});
\item \code{bool overflow() const} --- \code{flagbits \& kFlgOverflow}.
\end{itemize}
These are unsigned and are only meaningful when the ancestor list is filled
and long enough.  Because they are compiled methods they can be used in
\code{TTree::Draw()} / \code{Scan()} expressions once \code{libdk2nuTree}
is loaded, e.g.\ \code{"ancestor[indxp()].pdg"}.

\subsection{\code{bsim::Decay}}\label{sec:decay}

The core information about the neutrino and the decay that produced it.
Members marked \dag{} are required by \code{calcEnuWgt}
(\sect{calcenuwgt}) to compute the neutrino energy and weight at an
arbitrary point; files that omit them cannot be used with \genie.

\begin{fieldtable}{Decay}
norig     & Int\_t & -- & neutrino origin as recorded by the legacy
             simulations (\code{g4numi}: 1 = from a particle leaving the
             target or baffle, 2 = from scraping, 3 = from a muon decay; not
             filled by \code{flugg}). \open{current meaning, if any} \\
ndecay    & Int\_t & -- & decay process, \code{bsim::dkproc\_t}
             (\tab{dkproc}).  Used to select the muon-polarisation
             correction. \\
ntype\,\dag & Int\_t & -- & neutrino flavour (\textsc{pdg}). \\
vx, vy, vz\,\dag & Double\_t & cm & neutrino production (parent decay)
             vertex. \\
pdpx, pdpy, pdpz\,\dag & Double\_t & GeV/$c$ & momentum of the neutrino
             parent at the decay vertex. \\
ppdxdz, ppdydz\,\dag & Double\_t & -- & direction slopes of the neutrino
             parent at its production point. \\
pppz\,\dag & Double\_t & GeV/$c$ & $p_z$ of the neutrino parent at its
             production point (so $p_x = $ \code{ppdxdz*pppz}, etc.). \\
ppenergy\,\dag & Double\_t & GeV & energy of the neutrino parent at its
             production point. \\
ppmedium  & Int\_t & -- & code of the material in which the neutrino
             parent was produced (legacy \code{gnumi} codes,
             \tab{materials}). \\
ptype\,\dag & Int\_t & -- & neutrino parent species (\textsc{pdg}). \\
muparpx, muparpy, muparpz\,\dag & Double\_t & GeV/$c$ & momentum of the
             neutrino \emph{grandparent}.  If the parent (\code{ptype}) is a
             muon this is the grandparent at its decay point (i.e.\ where the
             muon was produced); otherwise it is the grandparent hadron at
             its production point.  Redundant with the ancestor list when
             that is filled. \\
mupare\,\dag & Double\_t & GeV & energy of the grandparent, as above. \\
necm\,\dag  & Double\_t & GeV & neutrino energy in the parent rest
             frame. \\
nimpwt\,\dag & Double\_t & -- & cumulative importance weight of the
             chain from production to this decay. \\
sumnimpwt2 & Double\_t & -- & cumulative sum of squares of the importance
             weights, needed for the statistical error of the
             importance-weighting process (added in \code{DK2NUVER} 11). \\
\end{fieldtable}

For a muon parent, \code{ppdxdz}, \code{ppdydz}, \code{pppz},
\code{ppenergy} describe the \emph{muon} at its production point and
\code{muparp*}/\code{mupare} describe the muon's parent (the pion or kaon)
at the point where it decayed; together they give the muon spin direction
used in the polarisation correction.

\begin{table}[htb]
\centering\small
\begin{tabular}{rll}
\toprule
\code{ndecay} & \code{bsim::dkproc\_t} & process \\ \midrule
 0 & \code{dkp\_unknown}         & unknown \\
 1 & \code{dkp\_k0l\_nuepimep}    & $K^0_L \to \nu_e\,\pi^-\,e^+$ \\
 2 & \code{dkp\_k0l\_nuebpipem}   & $K^0_L \to \bar\nu_e\,\pi^+\,e^-$ \\
 3 & \code{dkp\_k0l\_numupimmup}  & $K^0_L \to \nu_\mu\,\pi^-\,\mu^+$ \\
 4 & \code{dkp\_k0l\_numubpipmum} & $K^0_L \to \bar\nu_\mu\,\pi^+\,\mu^-$ \\
 5 & \code{dkp\_kp\_numumup}      & $K^+ \to \nu_\mu\,\mu^+$ \\
 6 & \code{dkp\_kp\_nuepi0ep}     & $K^+ \to \nu_e\,\pi^0\,e^+$ \\
 7 & \code{dkp\_kp\_numupi0mup}   & $K^+ \to \nu_\mu\,\pi^0\,\mu^+$ \\
 8 & \code{dkp\_kp\_numubmum}     & $K^- \to \bar\nu_\mu\,\mu^-$ \\
 9 & \code{dkp\_kp\_nuebpi0em}    & $K^- \to \bar\nu_e\,\pi^0\,e^-$ \\
10 & \code{dkp\_kp\_numubpi0mum}  & $K^- \to \bar\nu_\mu\,\pi^0\,\mu^-$ \\
11 & \code{dkp\_mup\_nusep}       & $\mu^+ \to \bar\nu_\mu\,\nu_e\,e^+$ \\
12 & \code{dkp\_mum\_nusep}       & $\mu^- \to \nu_\mu\,\bar\nu_e\,e^-$ \\
13 & \code{dk\_pip\_numumup}      & $\pi^+ \to \nu_\mu\,\mu^+$ \\
14 & \code{dk\_pim\_numubmum}     & $\pi^- \to \bar\nu_\mu\,\mu^-$ \\
15 & \code{dk\_mum\_capture}      & $\mu^- + (A,Z) \to \nu_\mu + \ldots$
                                    (capture; kinematics differ from 12) \\
16 & \code{dkp\_maximum}          & one past the last defined code \\
999 & \code{dkp\_other}           & other / unusual \\
\bottomrule
\end{tabular}
\caption{Decay process codes stored in \code{Decay::ndecay}
         (\code{bsim::dkproc\_t} in \code{dk2nu.h}).  Note that the
         enumerator names for codes 8--10 say \code{kp} although the
         parent is a $K^-$.}
\label{tab:dkproc}
\end{table}

\subsection{\code{bsim::NuRay}}\label{sec:nuray}

Each \code{NuRay} is the neutrino from this decay forced through one fixed
point.  \code{nuray[i]} corresponds to \code{DkMeta::location[i]}.

\begin{fieldtable}{NuRay}
px, py, pz & Double\_t & GeV/$c$ & neutrino momentum towards the location. \\
E          & Double\_t & GeV & neutrino energy at the location. \\
wgt        & Double\_t & -- & location weight (probability for the
             neutrino to pass through a 100\,cm radius disk at the
             location, see \sect{calcenuwgt}); does \emph{not} include
             \code{nimpwt}. \\
\end{fieldtable}

By convention \code{location[0]} is named \code{"random decay"} and
\code{nuray[0]} holds the neutrino direction actually chosen by the
beamline simulation for the decay, with \code{wgt}$=1$.  The remaining
entries are computed by \code{calcLocationWeights()}.  A flux spectrum at
location $i$ is therefore
\begin{lstlisting}[style=plain]
dk2nuTree->Draw("nuray[i].E", "decay.nimpwt*nuray[i].wgt");
\end{lstlisting}
normalised by the POT in the metadata.  These pre-computed rays are useful
for quick plots; a detector of any appreciable size sees a range of angles
and therefore of spectra, which is why \genie recomputes the ray for a
random point on a flux window (\sect{genie}).

\subsection{\code{bsim::Ancestor}}\label{sec:ancestor}

The ancestor list records the chain of particles from the primary proton
(\code{ancestor[0]}) to the neutrino (\code{ancestor[indxnu()]}, the last
element).  The neutrino's parent is \code{ancestor[indxp()]} and its
grandparent \code{ancestor[indxgp()]}.  The list may contain entries for
the same particle at successive interactions (e.g.\ an elastic or
quasi-elastic scatter), so the parent's \emph{production} is not
necessarily \code{ancestor[indxp()]}; see \code{parIndex}.

\begin{fieldtable}{Ancestor}
pdg       & Int\_t & -- & species (\textsc{pdg}).  Cleared to 0. \\
startx, starty, startz & Double\_t & cm & position where this particle
            starts (is created or scatters). \\
startt    & Double\_t & ns & time at the start point, relative to the
            primary proton.  Cleared to 0 (used by \genie for ray timing,
            \sect{genie-time}). \\
startpx, startpy, startpz & Double\_t & GeV/$c$ & momentum at the start
            point. \\
stoppx, stoppy, stoppz & Double\_t & GeV/$c$ & momentum at the end of this
            segment (where it interacts or decays). \\
polx, poly, polz & Double\_t & -- & polarisation vector. \\
pprodpx, pprodpy, pprodpz & Double\_t & GeV/$c$ & momentum of the
            \emph{parent} of this particle at the point where this particle
            was produced.  Compiled only with \code{KEEP\_ANCESTOR\_PPRODPXYZ}
            (currently defined).  \open{exact intended meaning and whether
            to retire in favour of \code{parIndex}} \\
nucleus   & Int\_t & -- & \textsc{pdg} code of the nucleus involved in the
            interaction that produced this particle.  Cleared to 0. \\
parIndex  & Int\_t & -- & index identifying this particle's parent in the
            chain.  The header comment reads ``from nu (0), parent (1)
            \ldots to proton (n)'', i.e.\ counted in the opposite order to
            the vector index.  Cleared to $-1$.  Of the converters only
            \code{convert\_flugg.C} fills it (copied from the \code{flugg}
            ntuple); the g4numi (MINERvA) and BooNE files examined leave it
            at $-1$ for every ancestor, and there each element's parent is
            simply the previous element.
            \open{confirm ordering convention} \\
proc      & string & -- & name of the process that created this particle.
            Cleared to \code{"<<no-process>>"}. \\
ivol      & string & -- & name of the volume where the particle starts.
            Cleared to \code{"<<no-volume>>"}. \\
imat      & string & -- & name of the material where the particle starts.
            Cleared to \code{"<<no-material>>"}. \\
\end{fieldtable}

\noindent Methods: \code{SetStartXYZT()}, \code{SetStartP()},
\code{SetStopP()}, \code{SetPProdP()} set the triplets;
\code{r()} $=\sqrt{x^2+y^2}$ at the start; \code{startpt()},
\code{startp()}, \code{stoppt()}, \code{stopp()}, \code{pprodpt()},
\code{pprodp()} give transverse and total momenta;
\code{IsInTarget()} returns true if \code{ivol} is one of a hard-coded
list of NuMI target volume names (\code{TargetM}, \code{Budal\_H},
\code{TGT1}, \code{TGT1001}--\code{TGT10xx}, \ldots).

Files converted from the \textsc{miner}$\nu$\textsc{a} \code{g4numi}
format were limited to 10 ancestors; when that limit was hit the
\code{kFlgOverflow} bit is set and the last entry is not the neutrino.  The
overflow flag has been found to be unreliable in some files, so code
should not depend on it alone.

\subsection{\code{bsim::TgtExit}}\label{sec:tgtexit}

Information about the neutrino's ancestor as it left the target.  This is
possibly, but not necessarily, the neutrino's parent.  It overlaps with the
ancestor list but is kept because it is filled by simulations that do not
record full ancestry.

\begin{fieldtable}{TgtExit}
tvx, tvy, tvz & Double\_t & cm & position where the ancestor exits the
           target. \\
tpx, tpy, tpz & Double\_t & GeV/$c$ & momentum of the ancestor as it exits
           the target. \\
tptype     & Int\_t & -- & species of the ancestor exiting the target
           (\textsc{pdg}). \\
tgen       & Int\_t & -- & generation of the neutrino parent in the cascade:
           1 = primary proton, 2 = particles produced by the proton
           interaction, 3 = their daughters, \ldots \\
\end{fieldtable}

\subsection{\code{bsim::Traj}}\label{sec:traj}

Optional track points, intended for crude visualisation of the history.

\begin{fieldtable}{Traj}
trkx, trky, trkz & Double\_t & cm & position of a recorded track point. \\
trkpx, trkpy, trkpz & Double\_t & GeV/$c$ & momentum at that point. \\
\end{fieldtable}
\open{which particle and which boundaries trigger a point; whether any
current simulation fills this}

\subsection{Flags and user extensions}\label{sec:extend}

\paragraph{\code{flagbits}.}  Bits are defined by \code{bsim::flgbitval\_t}:
\begin{center}\small
\begin{tabular}{lll}
\toprule
\code{kFlgOverflow}  & \code{0x00000001} & ancestor list was truncated \\
\code{kMaskReserved} & \code{0x0000FFFF} & bits reserved for the package \\
\code{kMaskUser}     & \code{0xFFFF0000} & bits available to users \\
\bottomrule
\end{tabular}
\end{center}
No other bits are currently assigned.  (The proposal suggested bits to flag
non-standard units or \geant particle codes; that was not adopted: units are
always as in \sect{conventions} and codes are always \textsc{pdg}.)

\paragraph{\code{vint}, \code{vdbl}.}  These vectors let a simulation store
additional per-entry information, for short-term studies, without changing
the format.  When every entry uses the same fixed mapping, the names of the
elements should be recorded in \code{DkMeta::vintnames} and
\code{DkMeta::vdblnames}.  If sizes vary from entry to entry the
interpretation is up to the user.

\paragraph{Additional branches.}  A simulation can also add its own
branches to \code{dk2nuTree} alongside \code{dk2nu}; readers that only set
the \code{dk2nu} branch address are unaffected.  The example
\code{test\_fill\_dk2nu.C} (\app{listings}) demonstrates this with a
\code{NonStd} class.

%===========================================================================
\section{The metadata tree: \code{bsim::DkMeta}}\label{sec:dkmeta}

{\setlength{\tabcolsep}{3pt}%
\begin{fieldtable}[G V U D]{DkMeta}
job        & Int\_t & -- & job number; links to \code{Dk2Nu::job}. \\
pots       & Double\_t & -- & protons-on-target represented by this job's
             entries.  Cleared to 0. \\
beamsim    & string & -- & program and version that generated the file
             (e.g.\ \code{"g4numi/v6r1"}). \\
physics    & string & -- & physics generator (e.g.\ \code{"fluka08"},
             \code{"g4.10.4"}). \\
physcuts   & string & -- & tracking cuts (e.g.\ \code{"threshold=0.1GeV"}). \\
tgtcfg     & string & -- & target configuration
             (e.g.\ \code{"minos/epoch3/-10cm"}). \\
horncfg    & string & -- & horn configuration
             (e.g.\ \code{"FHC/185A/LE/h1xoff=1mm"}). \\
dkvolcfg   & string & -- & decay volume configuration
             (e.g.\ \code{"helium"}, \code{"vacuum"}). \\
beam0x, beam0y & Double\_t & cm & proton beam centre at the start. \\
beam0z     & Double\_t & cm & $z$ where the proton beam starts. \\
beamhwidth, beamvwidth & Double\_t & cm & horizontal and vertical widths of
             the proton beam. \\
beamdxdz, beamdydz & Double\_t & -- & proton beam slopes. \\
location   & vector<\hspace{0pt}Location> & -- & locations for which \code{nuray}
             entries were computed. \\
vintnames  & vector<\hspace{0pt}string> & -- & names for the elements of
             \code{Dk2Nu::vint}. \\
vdblnames  & vector<\hspace{0pt}string> & -- & names for the elements of
             \code{Dk2Nu::vdbl}. \\
\end{fieldtable}}

\noindent The configuration strings are free-form; no standard syntax has
been imposed.

\subsection{\code{bsim::Location}}

\begin{fieldtable}{Location}
x, y, z & Double\_t & cm & position in beam coordinates. \\
name    & string & -- & identifying name, e.g.\ \code{"NOvA NearDet"}. \\
\end{fieldtable}

%===========================================================================
\section{The chosen ray: \code{bsim::NuChoice}}\label{sec:nuchoice}

\code{NuChoice} is not stored in a \dktwonu file.  It describes the neutrino
ray actually chosen by a consumer such as \code{GDk2NuFlux}, in both the
beam and the user (detector) coordinate systems, and is typically written
alongside the \code{Dk2Nu} entry in the event-generator output
(\sect{genie-access}).

{\setlength{\tabcolsep}{3pt}%
\begin{fieldtable}[G V U D]{NuChoice}
pdgNu      & int & -- & neutrino flavour (\textsc{pdg}). \\
xyWgt      & double & -- & location (``x-y'') weight for the chosen point,
             from \code{calcEnuWgt}. \\
impWgt     & double & -- & importance weight (copied from
             \code{Decay::nimpwt}). \\
impWgtSum2 & double & -- & copied from \code{Decay::sumnimpwt2}. \\
p4NuBeam   & TLorentzVector & GeV & 4-momentum in beam coordinates. \\
x4NuBeam   & TLorentzVector & cm, ns & 4-position in beam coordinates. \\
p4NuUser   & TLorentzVector & GeV & 4-momentum in user coordinates. \\
x4NuUser   & TLorentzVector & user & 4-position in user coordinates (by
             default m and s; see \sect{genie-coords}). \\
\end{fieldtable}}

\noindent The \genie weight returned by \code{GDk2NuFlux::Weight()} for
weighted generation is \code{impWgt*xyWgt} (times a window-tilt factor,
\sect{genie-gen}).

%===========================================================================
\section{Helper functions}\label{sec:helpers}

\subsection{Reading locations: \code{readWeightLocations}}\label{sec:readloc}

Detector positions are not hard-coded in simulation code; they are read
from a text file (\code{tree/readWeightLocations.h}):
\begin{lstlisting}[style=cxx,numbers=none]
void bsim::readWeightLocations(std::string locfilename,
                               std::vector<bsim::Location>& locations);
void bsim::readWeightLocations(std::string locfilename, bsim::DkMeta* dkmeta);
void bsim::printWeightLocations(std::vector<bsim::Location>& locations);
void bsim::printWeightLocations(bsim::DkMeta* dkmeta);
\end{lstlisting}
The file name may contain environment variables (expanded with
\code{gSystem->ExpandPathName}).  Lines beginning with \code{\#} are
comments; every other line is \code{x y z name}, positions in cm in beam
coordinates, the rest of the line (trimmed of blanks and quotes) being the
name.  The \code{DkMeta} variant first appends a location
\code{(0,0,0,"random decay")} and then appends the file's locations, so
that \code{location[0]} is always the random decay.

Three location files are provided in \code{etc/}:
\code{locations.txt} (NuMI and Booster detectors: MINOS Near/Far, NOvA
Near/Far, \textsc{miner}$\nu$\textsc{a}, MiniBooNE, MicroBooNE, SciBooNE;
listed in \app{listings}), \code{NuMIlocations.txt} and
\code{LBNElocations.txt}.  The comments in these files record the source
of each position.

\subsection{Computing rays: \code{calcLocationWeights}}

\begin{lstlisting}[style=cxx,numbers=none]
void bsim::calcLocationWeights(const bsim::DkMeta* dkmeta, bsim::Dk2Nu* dk2nu);
\end{lstlisting}
For every \code{dkmeta->location[i]} other than \code{"random decay"} this
calls \code{calcEnuWgt} and appends a \code{NuRay} (momentum along the
line from the decay vertex to the location, energy and weight) to
\code{dk2nu->nuray}.  It expects \code{"random decay"} to be location 0 and
\code{nuray[0]} to have already been filled by the caller; it prints an
error and returns early otherwise.  The \code{decay} members marked \dag{}
in \sect{decay} must be filled first.

\subsection{The energy/weight calculation: \code{calcEnuWgt}}\label{sec:calcenuwgt}

\begin{lstlisting}[style=cxx,numbers=none]
int bsim::calcEnuWgt(const bsim::Decay& decay, const TVector3& xyz,
                     double& enu, double& wgt_xy);
int bsim::calcEnuWgt(const bsim::Dk2Nu* dk2nu, const TVector3& xyz,
                     double& enu, double& wgt_xy);
\end{lstlisting}
Given a decay and a point \code{xyz} (beam coordinates, cm) this returns
the energy \code{enu} (GeV) of a neutrino from that decay travelling
towards the point, and the probability \code{wgt\_xy} that it passes
through a disk of radius $R = 100$\,cm centred on the point and facing the
decay.  It is a C++ transliteration of R.~Milburn's weighting routine
(NuMI-NOTE-BEAM-0109, MINOS-DocDB-109) as used in \code{gnumi}.

Let $m$ be the parent mass (chosen from \code{ptype}: $\pi^\pm$, $K^\pm$,
$K^0_{L,S}$/$K^0$, $\mu^\pm$, $\Omega^\pm$, $n/\bar n$), $\vec p$ the
parent momentum at decay (\code{pdp}), $\gamma$ and $\beta$ its boost,
$E^*$ the rest-frame neutrino energy (\code{necm}), $\vec d$ the vector
from the decay vertex to the point and $\theta$ the angle between $\vec p$
and $\vec d$.  Then
\begin{align}
  E_\nu &= \frac{E^*}{\gamma\,(1-\beta\cos\theta)}, \\
  w_{xy} &= \frac{1-\cos\left(\arctan(R/|\vec d|)\right)}{2}
           \;\left[\frac{1}{\gamma(1-\beta\cos\theta)}\right]^2 ,
\end{align}
i.e.\ the solid-angle fraction of the disk in the lab times the Jacobian of
the boost.  (Older versions used the small-angle form
$R^2/(4 d_z^2)$ for the solid angle.)  For a stopped parent
($|\vec p| = 0$) the boost factor is 1.

For muon decays (\code{ndecay} 11 or 12 only; not for muon capture, code
15) the weight is multiplied by a polarisation correction: the neutrino is
boosted to the muon rest frame, the muon's parent (\code{mupar*}) is
boosted to the muon production rest frame (using \code{ppenergy},
\code{ppdxdz}, \code{ppdydz}, \code{pppz}) to give the spin direction, and
with $\cos\theta_s$ the angle between them
\begin{align}
\nu_e,\bar\nu_e:&\quad 1-\cos\theta_s, \\
\nu_\mu,\bar\nu_\mu:&\quad
   \frac{(3-2x) - (1-2x)\cos\theta_s}{3-2x}, \qquad x = 2E^*/m_\mu .
\end{align}

Masses are the \geant{}4 10.3 values (since v01\_08); compiling with
\code{-DHISTORIC\_MASS} restores the older values used to reproduce legacy
files exactly.

Return codes: 0 success; 1 unknown parent type; 2 a muon decay producing
a neutrino that is neither $\nu_e$ nor $\nu_\mu$ type; 3 muon decay with
missing parent information; 4 negative polarisation weight.  On any error
\code{enu} and \code{wgt\_xy} are set to 0.

\paragraph{Normalisation.}  Because $w_{xy}$ is a probability for a disk of
area $\pi R^2 = \pi\times 10^4\,\mathrm{cm}^2$, the flux per unit area at a
point from $N_\mathrm{POT}$ protons is
\begin{equation}
  \frac{d\Phi}{dA} = \frac{1}{\pi R^2\,N_\mathrm{POT}}
      \sum_\text{entries} \code{nimpwt}\times w_{xy}.
\end{equation}

%===========================================================================
\section{Writing \dktwonu files}\label{sec:write}

A beamline simulation (or a converter) writes a \dktwonu file as follows;
\path{scripts/examples/test_fill_dk2nu.C} (\app{listings}) is a complete
example.
\begin{enumerate}
\item Create one \code{bsim::Dk2Nu} and one \code{bsim::DkMeta}; fill the
      locations with \code{readWeightLocations()}.
\item Create \code{dk2nuTree} with branch \code{"dk2nu"} of class
      \code{"bsim::Dk2Nu"} and \code{dkmetaTree} with branch
      \code{"dkmeta"} of class \code{"bsim::DkMeta"} (split level 99).
\item For each neutrino-producing decay:
  \begin{enumerate}[nosep]
  \item call \code{dk2nu->clear()} (essential: otherwise vectors keep
        entries from the previous decay);
  \item fill \code{job}, \code{potnum}, \code{jobindx}, \code{decay},
        \code{ppv*}, \code{tgtexit}, the \code{ancestor} list (proton
        first, neutrino last) and anything optional;
  \item push the simulation's own random-direction neutrino into
        \code{nuray[0]} with \code{wgt=1};
  \item call \code{calcLocationWeights(dkmeta, dk2nu)};
  \item \code{dk2nuTree->Fill()}.
  \end{enumerate}
\item At the end, fill \code{dkmeta->job}, \code{pots} and the
      descriptive strings, call \code{dkmetaTree->Fill()} once, and write
      both trees.
\end{enumerate}

%===========================================================================
\section{Reading \dktwonu files}\label{sec:read}

\subsection{Interactive \ROOT}

\code{scripts/load\_dk2nu.C} loads \code{libdk2nuTree} (and optionally
\code{libdk2nuGenie}) from \code{\$DK2NU/lib} and adds the include paths
needed by ACLiC:
\begin{lstlisting}[style=plain]
root $DK2NU/scripts/load_dk2nu.C myfile.root
root [] dkmetaTree->Scan("job:pots:location.name","","colsize=20")
root [] dk2nuTree->Draw("nuray[3].E","decay.nimpwt*nuray[3].wgt*(decay.ntype==14)")
root [] dk2nuTree->Draw("@ancestor.size()")
root [] dk2nuTree->Scan("decay.ntype:ancestor[indxnu()].pdg:ancestor[indxp()].pdg")
\end{lstlisting}
Its signature is
\begin{lstlisting}[style=cxx,numbers=none]
int load_dk2nu(bool genie_too = false, bool cnvtinc = false, bool verbose = false)
\end{lstlisting}
\code{genie\_too} also loads \code{libdk2nuGenie}; \code{cnvtinc} adds the
\code{scripts/convert/*} directories to the include path.  The \code{@} in
\code{@ancestor.size()} applies \code{size()} to the collection rather than
to each element.  The index into \code{nuray} depends on the location file
used to make the file; check \code{location.name} first.

\subsection{Compiled C++}

Set branch addresses on chains of both trees:
\begin{lstlisting}[style=cxx,numbers=none]
TChain* cflux = new TChain("dk2nuTree");   cflux->Add(pattern);
TChain* cmeta = new TChain("dkmetaTree");  cmeta->Add(pattern);
bsim::Dk2Nu*  dk2nu  = new bsim::Dk2Nu;
bsim::DkMeta* dkmeta = new bsim::DkMeta;
cflux->SetBranchAddress("dk2nu",&dk2nu);
cmeta->SetBranchAddress("dkmeta",&dkmeta);
\end{lstlisting}
See \path{scripts/examples/test_read_dk2nu.C} (\app{listings}).  Link
against \code{libdk2nuTree} and include \code{dk2nu/tree/dk2nu.h} etc.\
from the install area.

\subsection{Python: \code{pydk2nu}}\label{sec:python}

When built with \code{-DDK2NU\_PYTHON\_ENABLED=ON}, the module
\code{pydk2nu} exposes all the data classes (members read/write), the
\code{dkproc} and \code{flgbitval} enumerations, and a reader:
\begin{center}\small
\begin{tabularx}{\textwidth}{>{\ttfamily}l X}
\toprule
\textrm{\textbf{method}} & \textbf{description} \\ \midrule
dk2nuFileReader(files) & chain the given list of files (both trees). \\
decays()               & iterator over \code{Dk2Nu} entries. \\
metas()                & iterator over \code{DkMeta} entries. \\
first(), next()        & explicit iteration (\code{None} at end); also
                         \code{first\_meta()}, \code{next\_meta()}. \\
decay\_all\_through\_point(p) & loop in C++ over all entries; returns an
         $N\times 3$ \code{numpy} array of (\code{ntype}, $E_\nu$, weight)
         for point \code{p} (beam coordinates, cm). \\
decay\_all\_through\_points(ps) & as above for $M$ points in one pass;
         returns $N\times(1+2M)$. \\
Dk2Nu.decay\_through\_point(p) & (\,$E_\nu$, weight) for one entry. \\
\bottomrule
\end{tabularx}
\end{center}
The weight returned by the Python helpers already includes
\code{nimpwt} and is converted to a flux per cm$^2$:
$\code{nimpwt}\times w_{xy}\times 10^{-4}/\pi$, i.e.\ divided by the
$\pi\,(100\,\mathrm{cm})^2$ area of the disk used by \code{calcEnuWgt}.
Histogramming it and dividing by the summed \code{pots} gives
$\nu$/cm$^2$/POT.  \code{python/README.md}, \code{example\_off\_axis.py}
and \code{example\_decay\_parent.py} show complete examples.

%===========================================================================
\section{The \genie flux driver: \code{GDk2NuFlux}}\label{sec:genie}

\code{genie::flux::GDk2NuFlux} implements the \genie interfaces
\code{GFluxI} (flux driver), \code{GFluxExposureI} (exposure in POT) and
\code{GFluxFileConfigI} (file and XML configuration).  It is registered
with the \genie flux-driver factory under the name
\code{genie::flux::GDk2NuFlux}, so applications such as
\code{gevgen\_fnal} or experiment frameworks can instantiate it by name
once \code{libdk2nuGenie} is loaded.  For each neutrino requested it
\begin{enumerate}[nosep]
\item reads (or re-uses) a \code{Dk2Nu} entry;
\item picks a random point on a user-defined flux window;
\item computes the neutrino energy and weight for that point with
      \code{calcEnuWgt};
\item for unweighted generation, accepts or rejects the ray against a
      maximum weight, and keeps track of the equivalent POT;
\item transforms the ray it returns to the detector (``user'') coordinate
      system.
\end{enumerate}

\subsection{Loading files and configuration}\label{sec:genie-load}

\begin{lstlisting}[style=cxx,numbers=none]
virtual void LoadBeamSimData(const std::vector<std::string>& filenames,
                             const std::string& det_loc);
\end{lstlisting}
\code{det\_loc} is the name of a \code{<param\_set>} in the XML
configuration file (\sect{genie-xml}); it is loaded first, and the job
exits if it cannot be found.  Each entry in \code{filenames} is a path
whose \emph{file-name} part may be a wildcard (\code{TRegexp} wildcard
syntax: \code{*}, \code{?}, \code{[\ldots]}); directory parts are not
expanded.  A name without a wildcard is used as given.  A wildcard is
expanded by listing the directory, which for a remote (\code{root://})
directory requires the server to allow listing; if the directory cannot be
listed, an error is logged and the pattern contributes no files.  Names are
collected into a sorted set (so duplicates are dropped) and each file is
opened with \code{TFile::Open()}, so files may be streamed via xrootd.  A
file that cannot be opened is reported and skipped.  A file is added to the
chains only if it contains both \code{dk2nuTree} and \code{dkmetaTree}
(names changeable with \code{SetTreeNames()}).  The POT for each file is
the sum of \code{pots} over its \code{dkmeta} entries.

After loading, if no maximum weight has been set, the driver runs
\code{ScanForMaxWeight()} (\sect{genie-gen}), then picks a \emph{random
starting entry} in the chain, resets the exposure counters, and computes
the POT per entry (\sect{genie-pot}).

The XML file name defaults to \code{GNuMIFlux.xml} (the file shared with
\genie's \code{GNuMIFlux} driver; \code{doc/examples/GNuMIFlux.xml},
\sect{genie-nova}, is an example).  It can be changed with the
environment variable \code{GDK2NUFLUXXML} or with
\code{SetXMLFileBase()} (from \code{GFluxFileConfigI}).  The file is
located by \genie's standard search (\code{GXMLPATH}, then
\code{\$GENIE/config}).  Its root element must be \code{<gdk2nu\_config>}
or, for compatibility with \code{GNuMIFlux} files,
\code{<gnumi\_config>}.

\subsection{Coordinate systems, units and transformations}\label{sec:genie-coords}

Internally the driver works in beam coordinates with cm and ns, matching
the files.  The user (detector) system has its own orientation, origin and
units.  The user units are metres and seconds, the units \genie uses for
positions and times; this is the default and the only fully supported
choice.  The methods \code{SetLengthUnits()} and \code{SetTimeUnits()}
(and \code{<units>}, \code{<lunits>}, \code{<tunits>} in XML) exist, and
should be called before the transformation or window are defined, but
other units have not been fully implemented or tested (in particular the
ray timing, \sect{genie-time}, assumes metres and seconds).
The setters take the value of \code{genie::utils::units::UnitFromString()}
for the unit, e.g.\ \code{UnitFromString("m")}.  \code{LengthUnits()} and
\code{TimeUnits()} return the current user units relative to metres and
seconds (1 for the defaults, 0.01 for cm), as
\code{ROOTGeomAnalyzer::LengthUnits()} does; they are not the values passed
to the setters.

The transformation is
\begin{align}
  \vec x_\text{user} &= s\,R\,\vec x_\text{beam} + \vec x_0, &
  \vec p_\text{user} &= R\,\vec p_\text{beam},
\end{align}
where $R$ is the beam-to-user rotation (\code{SetBeamRotation()},
\code{<beamdir>}), $\vec x_0$ is the beam origin expressed in user
coordinates and user units (\code{SetBeamCenter()}, \code{<beampos>}),
and $s$ is the cm-to-user length scale.  Times are scaled by the
ns-to-user factor.  \code{Beam2User\{Pos,Dir,P4\}} and
\code{User2Beam\{Pos,Dir,P4\}} perform these transformations;
\code{GetBeamRotation()} and \code{GetBeamCenter()} return the settings.

\subsection{Flux window}\label{sec:genie-window}

\begin{lstlisting}[style=cxx,numbers=none]
void SetFluxWindow(TVector3 p1, TVector3 p2, TVector3 p3);   // user coords
void GetFluxWindow(TVector3& p1, TVector3& p2, TVector3& p3) const;
\end{lstlisting}
The three points (user coordinates and units) define a parallelogram:
\code{p1} is a corner and $\vec d_1 = $ \code{p2}$-$\code{p1},
$\vec d_2 = $ \code{p3}$-$\code{p1} are its edges, which should be
orthogonal (a warning is printed if they are not).  Each ray starts at a
point chosen uniformly on the window,
$\vec x = \vec p_1 + u_1\vec d_1 + u_2\vec d_2$.  The window should be
large enough, and far enough upstream, that every ray that could interact
in the geometry of interest crosses it.  A spherical ``window''
(\code{SetFluxSphere}) is declared but not yet supported: calling it is
fatal.

\subsection{Generating rays}\label{sec:genie-gen}

\paragraph{Weighted rays.}  For each ray the driver
\begin{enumerate}[nosep]
\item advances to the next entry, unless the current entry has been used
      fewer than \code{fNUse} times (\code{SetEntryReuse(n)},
      \code{<reuse>}); at the end of the chain it cycles back to the
      beginning (\code{SetNumOfCycles}, 0 = unlimited).  Normally only the
      \code{decay} part of the entry is read at this point
      (see \emph{Partial reading} below);
\item adds the POT equivalent of one ray to the exposure (\sect{genie-pot});
\item rejects the ray (without further work) if \code{ntype} is not in the
      list given to \code{SetFluxParticles()} (default
      \numu, \numubar, \nue, \nuebar).  Rejected flavours still count towards
      the exposure, so e.g.\ generating only intrinsic \nue events is
      correctly normalised;
\item picks a point on the window and calls \code{calcEnuWgt} to get
      $E_\nu$ and $w_{xy}$; the momentum is along the line from the decay
      vertex to that point;
\item sets the weight
      $w = \code{nimpwt}\times w_{xy}\times|\hat p\cdot\hat n|$,
      where the last factor (the projection onto the window normal) is
      applied unless \code{SetApplyWindowTiltWeight(false)};
\item once the ray is to be returned (every ray when generating weighted
      rays, otherwise only accepted rays), reads the rest of the entry,
      transforms the ray's 4-momentum and 4-position to user coordinates
      (the weight needs only beam coordinates), sets the times
      (\sect{genie-time}) and, if requested, moves the ray origin upstream
      (\code{upstreamz}).
\end{enumerate}
A ray with $E_\nu$ above the maximum energy (\code{SetMaxEnergy},
default 120\,GeV, possibly raised by the scan) is fatal, since that maximum
has already been used to normalise interaction probabilities.

\paragraph{Partial reading.}  Everything needed to weight a ray is in
\code{Dk2Nu::decay}.  If the \code{dk2nu} branch was written split (as the
standard writers do, with split level 99), the driver reads only the
\code{decay} sub-branch for each new entry, and the rest of the entry
(\code{ancestor}, \code{nuray}, \code{traj}, \ldots) only when a ray from
that entry is accepted.  The max-weight scan (\code{ScanForMaxWeight()})
therefore never reads full entries.  The check is made for each file in the
chain: entries of a file whose \code{dk2nu} branch is not split are read in
full, as before, so split and unsplit files can be mixed.  Because the
scan reads only \code{decay}, ROOT's \code{TTreeCache} would otherwise learn
to prefetch only those baskets.  The driver therefore registers every
branch with the cache once the scan is done, so generation makes the same
grouped reads as full reading.
\code{GetDk2Nu()} always returns a complete entry after
\code{GenerateNext()}.  \code{SetPartialRead(false)} or
\code{<partialread> false </partialread>} restores full reading of every
entry.

\paragraph{Unweighted rays.}  By default (\code{SetGenWeighted(false)})
\code{GenerateNext()} returns rays with weight 1: each weighted ray is
accepted with probability $w/w_\mathrm{max}$.  The maximum
$w_\mathrm{max}$ is estimated by \code{ScanForMaxWeight()}, which generates
\code{fMaxWgtEntries} weighted rays (default 2\,500\,000) and records the
largest weight and energy.  Then
\begin{itemize}[nosep]
\item $w_\mathrm{max} = f_w\times\max(w)$, with $f_w$ the weight fudge
      factor (default 1.05);
\item if the user floor (\code{SetMinMaxWeight()}, \code{<minmaxwgt>}) is
      larger, it is used instead;
\item if $f_E\times\max(E_\nu)$ exceeds the declared maximum energy, the
      maximum energy is raised to it ($f_E$ from \code{SetMaxEFudge()},
      default 1.05).
\end{itemize}
If a ray is later found with $w > w_\mathrm{max}$, the action depends on
\code{SetMaxWeightFailModel()} / \code{<maxwgtfail>}:
\begin{description}[nosep,font=\ttfamily]
\item[bump (0, default)] raise $w_\mathrm{max}$ to $f_w\,w$ and continue.
      Subsequent POT accounting uses the new maximum, so the exposure stays
      correct, but the acceptance efficiency drops.
\item[frozen (1)] keep $w_\mathrm{max}$; the ray is accepted, so such
      high-weight rays are under-represented.  This avoids changes of scale
      mid-job, at the cost of a small bias.
\item[abort (2)] print the configuration and abort.
\end{description}
\code{PrintConfig()} reports the scan result, the final maximum, and how
many times it was exceeded.  Choosing the scan length, fudge factors,
re-use count and floor is a trade-off between speed and bias; the example
configuration (\sect{genie-nova}) records tuned values for particular
file sets.

\subsection{POT accounting}\label{sec:genie-pot}

\code{GetTotalExposure()} and \code{UsedPOTs()} return the POT equivalent
of the rays drawn so far.  With $N_\mathrm{POT}$ the sum of \code{pots}
over all loaded files, $N_\mathrm{entries}$ the number of entries, $A$
the window area and $R = 100$\,cm,
\begin{equation}
  \mathrm{POT\ per\ ray} = \frac{\pi R^2}{A}\,
     \frac{N_\mathrm{POT}}{N_\mathrm{entries}}\,\frac{1}{w_\mathrm{max}},
\end{equation}
added for every ray drawn (including re-used entries and rejected
flavours), using the value of $w_\mathrm{max}$ at that moment.
\code{Clear()} resets the counters (the \genie event-generation driver calls
it with \code{"CycleHistory"}).  \code{NFluxNeutrinos()} and
\code{SumWeight()} report the number of rays drawn and their summed weight.

\subsection{Ray timing}\label{sec:genie-time}

If the ancestor list is filled, the time of the decay is taken from the
\code{startt} member of the neutrino's ancestor entry,
\code{ancestor[indxnu()]} (ns, relative to the primary proton).  The ray's time at the window is that time plus the light travel
time from the decay point to the window, stored in both \code{x4NuBeam}
(ns) and \code{x4NuUser} (user time units).  If \code{startt} is 0 the
driver warns (a few times) that the time is probably wrong.  When the ray
is moved to \code{upstreamz} (\code{MoveToZ0()}), its position is moved
back along the momentum to the given user $z$ and the time is shifted by
the corresponding travel time.  \code{GetDecayDist()} returns the
distance from the decay point to the current ray position in user units.

\subsection{Random numbers}\label{sec:genie-random}

The driver has no generator of its own and never sets a seed.  All random
numbers come from \genie's singleton
\code{genie::RandomGen::Instance()->RndFlux()}
(\code{Framework/Numerical/RandomGen.h}).  In \genie~v3 all the named
streams (\code{RndFlux}, \code{RndKine}, \code{RndGeom}, \ldots) return the
\emph{same} \code{TRandom3} (Mersenne Twister) object, so the flux driver,
geometry navigation and physics models share one sequence.

\paragraph{Seeding.}  \code{RandomGen::SetSeed(seed)} seeds that
\code{TRandom3}, ROOT's \code{gRandom}, and PYTHIA6/8 with the same value.
\genie applications call it from their \code{--seed} option; user code that
drives \genie directly must call it itself.  If no seed is ever set,
\code{RandomGen} is initialised with \code{kDefaultRandSeed} $= 65539$
(\code{Framework/Conventions/Controls.h}).  The old \code{\$GSEED}
environment variable is no longer read: if it is set, \genie exits with a
fatal error.

\paragraph{Where the driver draws.}
\begin{enumerate}[nosep]
\item \code{LoadBeamSimData()}: the starting entry,
      \code{RndFlux().Integer(N\_entries)}, so that jobs sharing the same
      files do not all begin with the same decays.
\item \code{GenerateNext\_weighted()}: two uniform numbers giving the
      fractions along the two window edge vectors (\sect{genie-window}),
      i.e.\ the point at which \code{calcEnuWgt} is evaluated.
\item \code{GenerateNext()}: one uniform number $r$ per weighted ray; the
      ray is accepted if $r < w/w_\mathrm{max}$ (\sect{genie-gen}).
\item \code{ScanForMaxWeight()}: called from \code{LoadBeamSimData()}, it
      generates \code{fMaxWgtEntries} weighted rays and so consumes two
      numbers for each ray of an accepted flavour before any event is
      produced.
\end{enumerate}

\paragraph{Consequences.}
\begin{itemize}[nosep]
\item The seed must be set \emph{before} the flux files are loaded.  The
      starting entry and the max-weight scan are drawn inside
      \code{LoadBeamSimData()}; if the seed is set afterwards they use the
      default seed, and every job starts at the same entry with the same
      scan result.
\item A run is reproducible for a fixed seed, configuration and software
      version.  Because the stream is shared, any change that alters how
      many numbers are consumed (physics model settings, geometry, the scan
      length or \code{<minmaxwgt>}, \code{<reuse>}, the flavour list) shifts
      every later number, so the flux rays cannot be held fixed while the
      physics is varied simply by keeping the seed.
\item Entries whose flavour is not in \code{SetFluxParticles()} are
      rejected before a window point is picked and before the acceptance
      step, so they consume no random numbers (but still count towards the
      exposure).
\end{itemize}

\subsection{XML configuration}\label{sec:genie-xml}

Each \code{<param\_set name="...">} groups the settings for one detector
location.  Elements are applied in the order given; set the units first and
\code{<beamdir>} before a six-value \code{<beampos>}.  Values may be
separated and grouped liberally with spaces, commas, semicolons, colons,
parentheses, brackets or \code{=}.

\begin{center}\small
\begin{longtable}{>{\ttfamily}p{0.20\textwidth} p{0.74\textwidth}}
\toprule
\textrm{\textbf{element}} & \textbf{meaning} \\ \midrule
\endhead
verbose & XML parsing verbosity (0--3). \\
using\_param\_set & apply another named \code{param\_set} first (recursive);
   later elements override it.  Fatal if not found. \\
units, lunits & user length unit, e.g.\ \code{m}, \code{cm}.  Set first. \\
tunits & user time unit, e.g.\ \code{s}, \code{ns}. \\
beamdir & beam-to-user rotation; \code{type} attribute:
   \code{"series"} --- a sequence of \code{<rotation axis="x|y|z"
   units="rad|deg">angle</rotation>} elements;
   \code{"thetaphi3"} --- six values, $(\theta,\phi)$ of the new $x$, $y$,
   $z$ axes \`a la \geant{}3, with a \code{units} attribute;
   \code{"newxyz"} --- nine values, three direction vectors \`a la
   \geant{}4 (each normalised).  Only the last \code{<beamdir>} is used;
   rotations are not cumulative across elements. \\
beampos & beam origin in user coordinates (3 values), or
   \code{( user xyz ) = ( beam xyz )} (6 values) stating that a beam-frame
   point (in \emph{user} length units) is at the given user point.  The
   six-value form needs \code{<beamdir>} first. \\
window & three \code{<point coord="det">x, y, z</point>} elements in user
   coordinates (\sect{genie-window}). \\
enumax & up to four values: maximum energy (GeV, default 120); energy
   fudge $f_E$ (1.05); weight fudge $f_w$ (1.05); number of rays in the
   max-weight scan (2\,500\,000). \\
minmaxwgt & floor for $w_\mathrm{max}$. \\
maxwgtfail & \code{bump} (0), \code{frozen} (1) or \code{abort} (2). \\
upstreamz & user $z$ to move ray origins to; if $|z| > 10^{30}$ rays stay
   on the window. \\
reuse & number of times each entry is used before moving on. \\
partialread & \code{true} (default) or \code{false} (also \code{yes}/%
   \code{no}, \code{on}/\code{off}, \code{1}/\code{0}): read only
   \code{decay} until a ray is accepted (\sect{genie-gen}). \\
\bottomrule
\end{longtable}
\end{center}
Unrecognised elements produce a warning and are ignored.

\subsection{Example configuration file}\label{sec:genie-nova}

\code{doc/examples/GNuMIFlux.xml} combines the configuration file
used by the NOvA experiment with \code{GDk2NuFlux}, the SBND and ICARUS
sets from the one used by SBND, and the sets from the one used by the DUNE
near detector (it predates the driver and so still uses
the \code{<gnumi\_config>} root element and \code{GNuMIFlux} header
comments, both of which are accepted).  It contains \code{param\_set}s for
MINOS, NOvA (NDOS, ND, FD and ``Rock'' variants, and as-built/historical
alignments), BNB, MicroBooNE, \textsc{miner}$\nu$\textsc{a}, ArgoNeuT,
SBND (with an ``SBND-Dirt'' variant), ICARUS, the DUNE 2x2 prototype with
\textsc{miner}$\nu$\textsc{a} in the NuMI beam (``ProtoDUNE-ND'',
``MiniRun4'', ``MiniRun5'' and variants) and the DUNE near detector in the
LBNF beam (``DUNEND'' and ``DUNENDROCK''), plus a DUNE far detector
set (``DUNE-FD'', formerly ``LBNE-FarDet\_01'').  The ``DUNE-FD'' set is
unverified: its baseline, beam angle, detector dimensions and window are
copied unchanged from the LBNE-era set and have not been checked against
the current DUNE far detector design.  The NOvA Near Detector set
illustrates a
\code{newxyz} rotation, a large window offset towards the beam, a high
re-use count, a weight floor, and the \code{frozen} failure mode:
\lstinputlisting[style=xml,firstline=223,lastline=285,firstnumber=223]%
  {examples/GNuMIFlux.xml}
\noindent and the ``rock'' variant inherits it, enlarging the window and
pushing the ray origin far upstream:
\lstinputlisting[style=xml,firstline=309,lastline=321,firstnumber=309]%
  {examples/GNuMIFlux.xml}
\noindent The Far Detector set instead uses the default \code{bump} model;
as its comment notes, with about one interaction per spill and POT counted
at the end, a change of $w_\mathrm{max}$ during the job is not a problem:
\lstinputlisting[style=xml,firstline=323,lastline=376,firstnumber=323]%
  {examples/GNuMIFlux.xml}
\noindent The MINOS Near Detector set (lines 54--125 of the file) shows
all three \code{<beamdir>} forms and both \code{<beampos>} forms for the
same transformation.

\subsection{Access to the current entry}\label{sec:genie-access}

\begin{itemize}[nosep]
\item \code{GetDk2Nu()}, \code{GetNuChoice()} --- the current entry and
      chosen ray;
\item \code{GetDkMeta()} --- the metadata entry whose \code{job} matches
      the current entry (looked up on demand);
\item \code{GetEntryNumber()}, \code{Index()} --- entry in the chain;
\item \code{GetBranchInfo()} --- tells the \genie output stage to add
      branches \code{dk2nu} (\code{bsim::Dk2Nu}) and \code{nuchoice}
      (\code{bsim::NuChoice}) to the event tree, so each generated event
      carries its flux history;
\item \code{GetMetaDataTree()} --- the \code{dkmeta} chain, for copying
      into the output;
\item \code{GetFileList()}, \code{PrintCurrent()}, \code{PrintConfig()}.
\end{itemize}

\subsection{Performance}

Event generation from \dktwonu files is slower than from pre-digested
\code{gsimple} files, because every ray requires reading an entry and most
rays are rejected by \genie's interaction-probability step.  A 2016 timing
study with the NOvA Near Detector geometry found that entry re-use
(\code{<reuse>} 10--50) gave the largest gain, and that files carrying
ancestor information were noticeably slower to read than files without.
Differences between beamline simulations in how importance weights are
assigned also matter.  The study is summarised in \code{dk2nu\_timing.pdf}
(R.~Hatcher, 2016).

Partial reading (\sect{genie-gen}) helps most when entries are re-used
little and few rays are accepted, because then most entries are read only
to be rejected.  With heavy re-use, entries are read rarely and the gain is
small.  The saving is mainly CPU time spent unpacking the \code{ancestor},
\code{nuray} and \code{traj} vectors, not bytes read.

%===========================================================================
\appendix

\section{Building and installing}\label{app:build}

The package is built with CMake (3.14 or later if building the Python
bindings).  \ROOT's \code{root-config} must be on the \code{PATH}
(with \code{ROOTSYS} set); for the \genie driver \code{genie-config} must
also be on the \code{PATH}.
\begin{lstlisting}[style=plain]
mkdir build && cd build
cmake -DCMAKE_INSTALL_PREFIX=/path/to/install [options] /path/to/dk2nu
make && make install
\end{lstlisting}

\begin{center}\small
\begin{tabularx}{\textwidth}{>{\ttfamily}l >{\ttfamily}l X}
\toprule
\textrm{\textbf{option}} & \textrm{\textbf{default}} & \textbf{effect} \\ \midrule
CMAKE\_BUILD\_TYPE & RelWithDebInfo & compiler flags: \code{RelWithDebInfo}
   is optimised with debug symbols (\code{-O2 -g}); \code{Release},
   \code{Debug} (unoptimised) and others as usual for CMake. \\
WITH\_GENIE & ON & build \code{libdk2nuGenie}. \\
GENIE\_ONLY & OFF & build only the \genie driver, against an installed
   tree library given by \code{DK2NUDATA\_DIR}. \\
DK2NUDATA\_DIR & "" & install prefix of the tree library (for
   \code{GENIE\_ONLY}). \\
WITH\_TBB & ON & link \ROOT's implicit multithreading library and TBB
   (\code{TBB\_LIB}). \\
COPY\_AUX & ON & install \code{etc/} and \code{scripts/}. \\
DK2NU\_PYTHON\_ENABLED & OFF & build \code{pydk2nu} (fetches
   \code{pybind11} from GitHub at configure time). \\
\bottomrule
\end{tabularx}
\end{center}

The \genie part uses the environment variables \code{GENIE} (source tree,
for headers under \code{\$GENIE/src}), \code{GENIE\_LIB},
\code{LIBXML2\_FQ\_DIR} (and \code{LIBXML2\_INC}), \code{LOG4CPP\_INC} and
\code{LOG4CPP\_LIB}.

The install area contains \code{lib/} (\code{libdk2nuTree},
\code{libdk2nuGenie} and their \code{\_rdict.pcm}/\code{.rootmap} files),
\code{include/dk2nu/tree/} and \code{include/dk2nu/genie/} (so code
includes \code{"dk2nu/tree/dk2nu.h"} with \code{-I<prefix>/include}, or
\code{"tree/dk2nu.h"} with \code{-I<prefix>/include/dk2nu}),
\code{etc/}, \code{scripts/}, and \code{python/<X.Y>/} for the module.
Set \code{DK2NU} to the install prefix for \code{load\_dk2nu.C} and the
example macros; add \code{\$DK2NU/lib} to the library path and
\code{\$DK2NU/python/<X.Y>} to \code{PYTHONPATH}.

%===========================================================================
\section{Converting legacy ntuples}\label{app:convert}

\code{scripts/convert/} contains \ROOT macros that convert the older
formats into \dktwonu files:

\begin{center}\small
\begin{tabular}{llll}
\toprule
\textbf{format} & \textbf{tree} & \textbf{macro} & \textbf{based on} \\ \midrule
\code{gnumi}     & \code{h10}    & ---                            & \geant{}3 \\
\code{flugg}     & \code{h10}    & \code{flugg/convert\_flugg.C}  & \textsc{fluka}+\geant{}4 \\
\code{g4numi}    & \code{nudata} & ---                            & \geant{}4 \\
\code{g4minerva} & \code{nudata} & \code{g4minerva/convert\_g4minerva.C} & \geant{}4 \\
\code{g4lbne}    & \code{nudata} & \code{g4lbne/convert\_g4lbne.C} & \geant{}4 \\
\bottomrule
\end{tabular}
\end{center}
Each format directory holds a \code{MakeClass}-generated reader
(\code{<fmt>.h}, \code{<fmt>.C}, produced by
\code{aux/mkgclasses3.sh}); \code{g3numi} and \code{g4numi} have readers
but no conversion macro.  \code{common\_convert.C} provides the shared
steps (\code{ConvertInit}, \code{ConvertBookNtuple},
\code{ConvertFinish}), \geant{}3$\to$\textsc{pdg} code translation
(\code{ConvertGeantToPdg}, \code{Convert5xToPdg}), POT estimation from
the highest proton number when the input lacks it, and cross-checks of the
recomputed energies and weights against the near/far values stored in the
input.  For example:
\begin{lstlisting}[style=plain]
root -b -q '$DK2NU/scripts/load_dk2nu.C(false,true)' \
  '$DK2NU/scripts/convert/g4lbne/convert_g4lbne.C+("in.root",42,"$DK2NU/etc/LBNElocations.txt")'
\end{lstlisting}
The arguments are input file, job number (the legacy files often did not
set a usable run number), a location file or cross-check selector (for
\code{flugg}: the FLUKA volumes file, then \code{"MINOS"} or \code{"NOvA"}
to choose which stored near/far positions to cross-check), and optional
maximum entries and dump frequency.  The output is
\code{<input>\_to\_dk2nu.root}.  The legacy element names map to \dktwonu
as follows (capitalisation varied between formats):

\begin{center}\small
\begin{tabular}{ll}
\toprule
\textbf{legacy} & \textbf{\dktwonu} \\ \midrule
\code{run}, \code{evtno} & \code{job}, \code{potnum} \\
\code{Ndxdz}, \code{Ndydz}, \code{Npz}, \code{Nenergy} & \code{nuray[0]} (random decay) \\
\code{NdxdzNear}, \code{NenergyN}, \code{NWtNear}, \ldots & recomputed into \code{nuray[i]} \\
\code{Norig}, \code{Ndecay}, \code{Ntype} & \code{decay.norig}, \code{decay.ndecay}, \code{decay.ntype} \\
\code{Vx}, \code{pdPx}, \code{ppdxdz}, \code{pppz}, \code{ppenergy}, \ldots & \code{decay.vx}, \code{decay.pdpx}, \ldots \\
\code{ppmedium}, \code{ptype}, \code{muparpx}, \ldots, \code{Necm}, \code{Nimpwt} & \code{decay.*} \\
\code{ppvx}, \code{ppvy}, \code{ppvz} & \code{ppvx}, \code{ppvy}, \code{ppvz} \\
\code{tvx}, \code{tpx}, \code{tptype}, \code{tgen}, \ldots & \code{tgtexit.*} \\
\code{pdg[]}, \code{start*[]}, \code{stopp*[]}, \code{pprodp*[]}, \code{proc[]}, \code{ivol[]} & \code{ancestor[]} (mm$\to$cm, MeV$\to$GeV) \\
\code{overflow} & \code{flagbits \& kFlgOverflow} \\
\code{trkx[]}, \code{trkpx[]}, \ldots & \code{traj[]} \\
\bottomrule
\end{tabular}
\end{center}
Legacy elements with no \dktwonu equivalent (\code{xpoint},
\code{tgptype}, \code{tgppx}, \code{tprivx}, \code{beamx}, \code{beampx},
\code{trackId}, \code{parentId}, \code{fvol}, the \code{flugg}
derived quantities \code{Vr}, \code{pdP}, \code{tp}, \ldots) are dropped;
per-file beam parameters belong in \code{DkMeta}.

%===========================================================================
\section{Material codes}\label{app:materials}

\begin{table}[htb]
\centering\small
\begin{tabular}{rl@{\qquad}rl}
\toprule
\textbf{code} & \textbf{material} & \textbf{code} & \textbf{material} \\ \midrule
 5 & Beryllium & 20 & Shotcrete \\
 6 & Carbon    & 21 & Variable-density aluminium \\
 9 & Aluminium & 22 & Variable-density steel \\
10 & Iron      & 23 & 1018 steel \\
11 & Slab steel & 24 & A500 steel \\
12 & Blu steel  & 25 & Water \\
15 & Air        & 26 & M1018 steel \\
16 & Vacuum     & 28 & Decay-pipe vacuum \\
17 & Concrete   & 31 & CT852 \\
18 & Target (C) & & \\
19 & Rebar concrete & & \\
\bottomrule
\end{tabular}
\caption{Material codes stored in \code{Decay::ppmedium}, as defined by
\code{gnumi} and used in the NuMI flux files.  Other simulations may use
their own numbering; \code{Ancestor::imat} records material names
directly.}
\label{tab:materials}
\end{table}

%---------------------------------------------------------------------------
\section{Code listings}\label{app:listings}

All listings are included directly from the repository.

\subsection{\code{tree/dk2nu.h}}
\lstinputlisting[style=cxx]{../tree/dk2nu.h}

\subsection{\code{tree/dkmeta.h}}
\lstinputlisting[style=cxx]{../tree/dkmeta.h}

\subsection{\code{tree/NuChoice.h}}
\lstinputlisting[style=cxx]{../tree/NuChoice.h}

\subsection{\code{tree/dflt.h}}
\lstinputlisting[style=cxx]{../tree/dflt.h}

\subsection{\code{tree/readWeightLocations.h}}
\lstinputlisting[style=cxx]{../tree/readWeightLocations.h}

\subsection{\code{tree/calcLocationWeights.h}}
\lstinputlisting[style=cxx]{../tree/calcLocationWeights.h}

\subsection{\code{tree/calcLocationWeights.cxx}}
\lstinputlisting[style=cxx]{../tree/calcLocationWeights.cxx}

\subsection{\code{genie/GDk2NuFlux.h}}
\lstinputlisting[style=cxx]{../genie/GDk2NuFlux.h}

\subsection{\code{etc/locations.txt}}
\lstinputlisting[style=plain]{../etc/locations.txt}

\subsection{\code{scripts/load\_dk2nu.C}}
\lstinputlisting[style=cxx]{../scripts/load_dk2nu.C}

\subsection{\code{scripts/examples/test\_fill\_dk2nu.C}}
\lstinputlisting[style=cxx]{../scripts/examples/test_fill_dk2nu.C}

\subsection{\code{scripts/examples/test\_read\_dk2nu.C}}
\lstinputlisting[style=cxx]{../scripts/examples/test_read_dk2nu.C}

\subsection{\code{scripts/examples/test\_read\_locations.C}}
\lstinputlisting[style=cxx]{../scripts/examples/test_read_locations.C}

%---------------------------------------------------------------------------
\section{Open questions and known issues}\label{app:issues}

Items marked \open{\ldots} in the text, collected here, together with
behaviour noticed while preparing this document that users should be aware
of.  None of these change the file format.

\begin{enumerate}
\item \code{Decay::norig}: is it still filled meaningfully by current
      simulations?
\item \code{Ancestor::pprodp\{x,y,z\}}: precise meaning, and whether to
      retire it (the header notes the intent to replace it with
      \code{parIndex}).
\item \code{Ancestor::parIndex}: the header comment counts from the neutrino
      (0) to the proton ($n$), opposite to the vector order.  This could not
      be checked against data: in 20\,000 entries from each of two g4numi
      (MINERvA) and two BooNE files, every ancestor has the cleared value
      $-1$.  Only \code{convert\_flugg.C} fills it, from the \code{flugg}
      ntuple, so confirming the convention needs a \code{flugg}-derived
      file or the code that fills that ntuple.
\item \code{Traj}: which particle and boundaries generate points; is it
      filled by any current simulation?
\item User units other than metres and seconds in \code{GDk2NuFlux} were
      never fully implemented or tested; for example the light-travel time
      added in \sect{genie-time} divides a distance in user length units by
      $c$ in m/s.  Possible future work.
\item Flux spheres (\code{SetFluxSphere()}) are not supported in ray
      generation; calling it is fatal.
\end{enumerate}

\end{document}
